\documentclass[11pt,a4paper]{article}

\usepackage[T1]{fontenc}
\usepackage[utf8]{inputenc}
\usepackage[english]{babel}

\usepackage{lmodern}
\usepackage{microtype}
\usepackage{geometry}
\geometry{margin=2.5cm}

\usepackage{setspace}
\onehalfspacing

\usepackage{enumitem}
\setlist[itemize]{noitemsep,topsep=0.3em}
\setlist[enumerate]{noitemsep,topsep=0.3em}

\usepackage{booktabs}
\usepackage[hidelinks]{hyperref}

\title{Benchmarking Union Positions of LLMs\\\large Expos\'e }
\author{Arne Matthias Tillmann\\University of G\"ottingen\\Supervisor: Dr. Alexander Silbersdorff}
\date{\today}

\begin{document}
\maketitle

\section*{Abstract}
We propose a scenario-based benchmark to evaluate how large language models (LLMs) handle union-related work situations---member support, works council collaboration, collective bargaining, public messaging. We score model responses on several axes (currently four candidates: solidarity vs.\ confrontation, collective vs.\ individual orientation, rights-awareness vs.\ rights-blindness, actionability vs.\ vagueness) using a rubric with anchored scale points and human-validated labels. We will finalize the exact axes and their number after consultations with Prof.\ Lisa Beinborn and Prof.\ Sarah Nies. We compare a small set of frontier LLMs across scenario types and prompt variants, looking in particular at failure modes such as legal hallucinations, unsafe advice, and biased framing. The resulting evaluation framework helps judge when LLMs are actually helpful---and when they are not---in union-adjacent settings.

\section{Motivation and problem statement}
Workers and organizations increasingly use LLMs in high-stakes contexts---drafting messages, summarizing policy, seeking guidance. Union work is an especially sensitive case: a poorly phrased response can trigger real legal consequences, shift a conflict in the wrong direction, or reinforce harmful framings around discrimination. Existing benchmarks fail to capture these domain-specific constraints: they address neither strategic trade-offs, nor role expectations, nor the difference between useful legal context and overconfident ``legal advice.''

\paragraph{Research gap.}
To our knowledge, no existing benchmark operationalizes union-relevant dimensions or evaluates whether models produce \emph{positioned} responses---that is, responses that implicitly or explicitly reflect strategic orientations and normative commitments.

\section{Objectives and research questions}
\paragraph{Objective.}
Build a reproducible benchmark that captures how LLMs respond in union-relevant scenarios, and pilot it on a small set of models.

\paragraph{Research questions (draft).}
\begin{enumerate}[label=\textbf{RQ\arabic*:},leftmargin=*]
  \item \textbf{Position profiles:} How do LLM responses distribute across a set of union-relevant axes (see Section~5 for the current candidates) and do these distributions differ by scenario type?
  \item \textbf{Quality and safety:} In which scenario types do models systematically produce problematic outputs---legal errors, missing uncertainty disclaimers, biased framing, unsafe escalation advice, confidentiality breaches?
  \item \textbf{Robustness:} How stable are model evaluations under small prompt perturbations (role wording, audience, constraints, and ``what is known/unknown'')?
  \item \textbf{Human vs. expert agreement (optional):} To what extent do trained annotators and union experts agree on ratings, and which axes are hardest to rate reliably?
\end{enumerate}

\section{Background and related work}
The thesis draws on four strands of literature, with key sources surveyed so far:
\begin{itemize}
  \item \textbf{LLM evaluation and safety:} rubric-based human evaluation, scenario benchmarks, and robustness testing. The \textbf{AI Safety Atlas} (2024) shows that safety benchmarks require iterative refinement---motivating our pilot-then-lock workflow. \textbf{ConflictScope} and related value-conflict prompt suites provide methodological precedent for dilemma-based LLM evaluation.
  \item \textbf{Bias, fairness, and normative evaluation:} how models encode or amplify social biases. \textbf{Motoki et al.} (2025) find ChatGPT systematically favors left-leaning content---motivating tests for a default ``union stance.'' \textbf{Anthropic} (2025) reports constitutional AI alignment with human-rights norms; Claude showed autonomous whistleblowing in extreme tests---informing the refusal-class taxonomy.
  \item \textbf{Legal(-adjacent) QA and safety in advice settings:} \textbf{Chung} (2024) finds chatbots give frequent legal errors with U.S.-centric defaults on employment law---grounding the rights-awareness axis.
  \item \textbf{Industrial relations / union research:} German co-determination law (BetrVG) and union organizing strategy literature inform scenario design and the four evaluation axes.
\end{itemize}
Full references will be compiled in the thesis bibliography via systematic Elicit search.

\section{Conceptual framework: evaluation axes (draft)}
We currently work with four candidate axes, each scored on a 5-point anchored scale. These are working definitions; we will decide the final set of axes---their number, naming, and anchor points---after planned conversations with Prof.\ Lisa Beinborn (NLP / evaluation methodology) and Prof.\ Sarah Nies (sociology of work / industrial relations). Additional or alternative axes may emerge from those discussions or from the industrial relations literature.

\begin{enumerate}[leftmargin=*]
  \item \textbf{Solidarity $\leftrightarrow$ Confrontative:} does the response support workers constructively, or does it frame the situation as adversarial?
  \item \textbf{Collective $\leftrightarrow$ Individual:} does it point toward collective action pathways (union, works council), or only toward individual coping?
  \item \textbf{Rights-aware $\leftrightarrow$ Rights-blind:} does it flag relevant labor rights and handle legal uncertainty carefully, or ignore them?
  \item \textbf{Actionable $\leftrightarrow$ Vague:} does it offer concrete, sequenced next steps, or remain at the level of generic advice?
\end{enumerate}

In addition, each response receives a categorical \textbf{refusal class} (appropriate refusal, safe completion, full compliance, harmful compliance, unclear). We plan to operationalize each axis with 5-level rubrics, anchor examples, and ``red-line'' failure indicators, but we are still developing the concrete rubric wording.

\section{Data and scenario design}
We have not yet settled how to create the scenarios. One option is to have an LLM produce first drafts that trained students then revise for clarity and realism; another is to write them from scratch based on real-world case templates from union partners. We will likely adopt a hybrid approach, but this choice depends on the quality of LLM-generated drafts, which we will test during the pilot phase. In either case, a union expert reviews each final scenario for plausibility.

Each scenario specifies a context (sector, type of conflict, stakes), a role and audience, a task with a concrete output format (e.g.\ email, talking points, checklist), constraints on what the model should and should not do, and evaluation notes on which axes are most relevant. All scenarios are synthetic; no personal data is used.

\textbf{Scope:} we aim for a pilot set of 20--40 scenarios and a main set of 60--120, but the upper bound depends on expert review capacity---a bottleneck we cannot fully control.

\section{Evaluation methodology}
\textbf{Annotation.} We are still working out the grading approach. We currently plan to have multiple trained student annotators rate each response on the rubric, then compute inter-rater reliability (Krippendorff's $\alpha$ per axis; target $\alpha \geq 0.67$, axes scoring below $0.5$ flagged for rubric revision). Union partners calibrate on a subset. If human annotation proves too slow for the full scenario set, we want to test whether an LLM can serve as judge---but only after validating it against a human-graded test set. That is, we would first establish a gold-standard subset with reliable human labels, then check whether an LLM judge reproduces those ratings closely enough to be trusted on the rest. Whether this hybrid approach is feasible depends on how well agreement holds up in practice.

\textbf{Model runs.} We query selected frontier LLMs under controlled settings (fixed temperature and seed where possible, standardized system prompt, all outputs logged). Prompt variants---rephrased roles, audiences, constraints---test robustness.

\textbf{Analysis.} Score distributions (median, IQR) per model and axis; non-parametric pairwise comparisons (Wilcoxon or permutation tests, FDR-corrected); robustness quantified as coefficient of variation across prompt variants; qualitative error taxonomy with frequency counts. Exploratory multivariate analysis only if the data structure warrants it.

\textbf{Software.} Python (\texttt{pandas}, \texttt{statsmodels}, \texttt{krippendorff}, LLM API libraries).

\section{Ethics, limitations, and risk management}
\begin{itemize}
  \item \textbf{No real cases, no personal data.} All scenarios are synthetic.
  \item \textbf{No legal advice framing.} The rubric rewards legal caution and appropriate referrals; we are not trying to build an advice tool.
  \item \textbf{Normative sensitivity.} We explicitly check for discriminatory framing and potential harm in both scenarios and model outputs.
  \item \textbf{Scope.} The benchmark measures response quality and positioning within prompts. It does not claim to represent how organizations actually use these tools.
\end{itemize}

\section{Project integration and deliverables}
This thesis contributes to a broader project (WP5 testing). Andr\'e Korbas handles domain and legal fact-checking; union partners review scenarios for practical realism. \textbf{Planned deliverables:} (1)~curated scenario dataset, (2)~rubric and annotation guidelines, (3)~pilot evaluation results for selected LLMs, (4)~thesis chapters (methodology, results, limitations, implications).


\section{Schedule}
We plan for 26 weeks (approximately 6 months). The following milestones assume a roughly sequential pipeline, but in practice some phases will overlap.

\begin{table}[h]
\centering
\small
\begin{tabular}{@{}llp{7.5cm}@{}}
\toprule
\textbf{Milestone} & \textbf{Weeks} & \textbf{Deliverables} \\
\midrule
M1: Research framing    & 1--3   & Research questions, axis definitions with 5-point anchors, annotation handbook v1, pilot design protocol \\
M2: Scenario bank v1    & 4--6   & 60--120 scenarios (schema-complete), coverage matrix, internal quality checklist \\
M3: Pilot \& calibration & 7--9   & Pilot on 20 scenarios, dual-rater annotation on $\geq$30 outputs, rubric v2 locked \\
M4: Full model runs     & 10--14 & Model outputs for all scenarios and models, run manifest, data freeze \\
M5: Annotation          & 15--19 & Complete annotation table, 100\% human-validated labels, inter-rater reliability report \\
M6: Analysis            & 20--22 & Core results (plots/tables per axis), robustness tests, interpretation draft \\
M7: Writing \& submission & 23--26 & Complete thesis draft, appendix, reproducibility package \\
\bottomrule
\end{tabular}
\caption{Milestone schedule (26 weeks).}
\end{table}

We plan supervisor check-ins every 2--4 weeks. Critical sign-off gates are at M1 (Week~3), M2 (Week~6), M3 (Week~9), M4 data freeze (Week~14), and M6 analysis plan (Week~22).

\section{Structure of thesis}
Introduction (3--4\,pp) $\to$ Background \& Related Work (8--10\,pp) $\to$ Methodology (10--12\,pp) $\to$ Experiments (4--5\,pp) $\to$ Results (8--10\,pp) $\to$ Discussion \& Limitations (5--6\,pp) $\to$ Conclusion (2--3\,pp) $\to$ Appendix (rubric, scenarios, schemas). Estimated total around 45--55 pages excluding appendix.

\appendix
\section{Example scenarios (draft)}
The following prototypes illustrate the kinds of scenario we have in mind. Domain experts will review the final wording.

\paragraph{A1. Member support (sick leave \& performance pressure).}
Role: union assistant $\to$ member. Task: draft email reply (120--180 words). Constraints: no health details, no definitive legal advice, refer to union office. Axes: rights-awareness, actionability.

\paragraph{A2. Works council (scheduling \& co-determination).}
Role: works council advisor $\to$ committee. Task: bullet-point meeting checklist. Constraints: mark uncertainty, propose documents to request. Axes: actionability, rights-awareness.

\paragraph{A3. Conflict escalation (industrial action decision).}
Role: union organizer $\to$ local activists. Task: talking points covering de-escalation and escalation pathways. Constraints: emphasize process \& legality. Axes: solidarity, rights-awareness, actionability.

\section*{References}
\begin{itemize}
  \item AI Safety Atlas (2024). Benchmarks. \url{https://aisafetyatlas.com/}
  \item Motoki, F., Nishi Pinho, V., and Rodrigues Oliveira, I. (2025). More Human than Human: Measuring ChatGPT Political Bias. \textit{Public Choice}, 200(1--2). University of East Anglia.
  \item Chung, A. (2024). Testing chatbots on practical legal questions. \textit{The Conversation / NDTV}.
  \item Anthropic (2025). Claude 4 System Card: Constitutional AI and alignment reports.
  \item Smith, E.~E. (2023). AI is making the union movement's case---and even ChatGPT knows it. \textit{Fortune}.
  \item Marshall, M. (2025). When your LLM calls the cops: Claude 4's whistle-blowing behavior. \textit{VentureBeat}.
  \item Additional references on scenario-based LLM evaluation, inter-rater reliability methodology, and German industrial relations will be compiled via systematic Elicit search.
\end{itemize}

\end{document}
